% !TEX root = ../matan.tex

\section{Функции нескольких вещественных переменных. Непрерывность и непрерывность вдоль направления.}

Пространство $\R^d$ снабжается евклидовой нормой
\[
|x-y|=\sqrt{\sum_{k=1}^{d}(x_k-y_k)^2},\qquad x,y\in\R^d.
\]
\emph{Областью} называется связное открытое множество $\Omega\subset\R^d$. Рассматриваются функции $f\colon\Omega\to\R$ (а в векторном случае --- отображения $f\colon\Omega\to\R^n$).

\begin{Definition}{Предел и непрерывность}{}
	Пусть $f\colon\Omega\subset\R^d\to\R$ и $x\in\Omega$ --- предельная точка множества $\Omega$. Говорят, что $f$ \emph{непрерывна} в точке $x$, если
	\[
	\lim_{\substack{y\to x\\ y\in\Omega}} f(y)=f(x),
	\]
	то есть для любого $\eps>0$ найдётся $\delta>0$ такое, что из $y\in\Omega$, $|y-x|<\delta$ следует $|f(y)-f(x)|<\eps$. Функция, непрерывная в каждой точке области, называется непрерывной на $\Omega$.
\end{Definition}

\begin{Definition}{Непрерывность вдоль направления}{}
	Пусть $x$ --- внутренняя точка $\Omega$ и $v\in\R^d$, $v\ne0$. \emph{Сужением $f$ на прямую, проходящую через $x$ в направлении $v$,} называется функция одной переменной
	\[
	\psi_v(t)=f(x+tv),\qquad t\in\{t\in\R\st x+tv\in\Omega\}.
	\]
	Функция $f$ \emph{непрерывна в точке $x$ вдоль направления $v$}, если $\psi_v$ непрерывна в нуле, то есть $\displaystyle\lim_{t\to0}f(x+tv)=f(x)$.

	В частности, для координатных направлений $v=e_k$ это означает непрерывность $f$ по $k$-й переменной при фиксированных остальных.
\end{Definition}

\begin{Statement}{}{}
	Если $f$ непрерывна в точке $x$, то для каждого $v\ne0$ функция $\psi_v(t)=f(x+tv)$ непрерывна в нуле.
\end{Statement}

\begin{proof}
	Зафиксируем $\eps>0$. По непрерывности $f$ в точке $x$ найдётся $\delta>0$ такое, что $|f(y)-f(x)|<\eps$ при $|y-x|<\delta$. Положим $y=x+tv$; тогда $|y-x|=|t|\,|v|$, и при $|t|<\delta/|v|$ имеем $|y-x|<\delta$, а значит
	\[
	|\psi_v(t)-\psi_v(0)|=|f(x+tv)-f(x)|<\eps.
	\]
	Это и есть непрерывность $\psi_v$ в нуле.
\end{proof}

\begin{Remark}{Обратное неверно}{}
	Непрерывность вдоль каждого направления (и тем более вдоль каждой координатной оси) \emph{не} влечёт непрерывности функции в точке. Причина в том, что приближение $y\to x$ в $\R^d$ при $d\ge2$ может происходить по произвольным --- в том числе криволинейным --- путям, а не только по прямым. Ниже два контрпримера нарастающей силы.
\end{Remark}

\begin{Example}{Непрерывность по осям, но разрыв в нуле}{}
	Рассмотрим
	\[
	f(x_1,x_2)=
	\begin{cases}
		\dfrac{x_1 x_2}{x_1^2+x_2^2}, & (x_1,x_2)\ne(0,0),\\[6pt]
		0, & (x_1,x_2)=(0,0).
	\end{cases}
	\]
	Вдоль координатных осей $f(t,0)=f(0,t)=0\to0=f(0,0)$, то есть $f$ непрерывна в нуле по каждой переменной в отдельности. Однако вдоль прямой $x_1=x_2=t$
	\[
	f(t,t)=\frac{t^2}{2t^2}=\frac12\xrightarrow{t\to0}\frac12\ne f(0,0),
	\]
	поэтому общего предела в нуле нет и $f$ разрывна в $(0,0)$.
\end{Example}

\begin{Example}{Непрерывность вдоль всех прямых, но разрыв в нуле}{}
	Рассмотрим
	\[
	f(x_1,x_2)=
	\begin{cases}
		\dfrac{x_1 x_2^2}{x_1^2+x_2^4}, & (x_1,x_2)\ne(0,0),\\[6pt]
		0, & (x_1,x_2)=(0,0).
	\end{cases}
	\]
	Для произвольного направления $v=(v_1,v_2)\ne0$ при $v_1\ne0$
	\[
	f(tv_1,tv_2)=\frac{t v_1\cdot t^2 v_2^2}{t^2 v_1^2+t^4 v_2^4}
	=\frac{t\,v_1 v_2^2}{v_1^2+t^2 v_2^4}\xrightarrow{t\to0}0,
	\]
	а при $v_1=0$ имеем $f(0,tv_2)\equiv0$. Значит, вдоль \emph{каждой} прямой предел в нуле равен $0=f(0,0)$, то есть $f$ непрерывна в нуле вдоль любого направления. Тем не менее вдоль параболы $x_1=x_2^2$ (в точках вида $(t^2,t)$):
	\[
	f(t^2,t)=\frac{t^2\cdot t^2}{t^4+t^4}=\frac{t^4}{2t^4}=\frac12\xrightarrow{t\to0}\frac12\ne0,
	\]
	поэтому общего предела в нуле не существует и $f$ разрывна в $(0,0)$.
\end{Example}
